\section{Reasoning and argumentation boundaries}
\label{sec:p18-r56-proofs}

\begin{annexitem}{annex:p18.r56.consequence-policy}{thm:p18-r56-consequence-policy}{A consequence relation does not select a repair policy}
\begin{proof}
For the fixed interpretation class, suppose every interpretation
satisfying all of \(\Gamma\) satisfies \(\varphi\).
If an interpretation satisfies all of \(\Delta\), with
\(\Gamma\subseteq\Delta\), it satisfies all of \(\Gamma\)
and hence \(\varphi\). This is exactly
\(\Delta\models\varphi\). No finiteness assumption is needed for
this clause. In particular adding premises within this same semantic
definition cannot by itself destroy a consequence; changing the
model class, inference policy or meaning of the premises is a
different operation.

For the finite classical set \(K\), all three members cannot be
true: \(p\) and \(p\to q\) force \(q\), contrary to \(\neg q\).
Each two-member subset is consistent, with these valuations:
\[
\begin{array}{c|cc}
\text{deleted member}&p&q\\\hline
p&0&0\\
p\to q&1&0\\
\neg q&1&1
\end{array}
\]
The first retains \(p\to q,\neg q\), the second \(p,\neg q\),
and the third \(p,p\to q\). A proper superset within \(K\)
must be the inconsistent full set, so each pair is maximal consistent.
Every singleton is contained in one of these consistent pairs and
the empty subset is contained in them all. There are therefore exactly
these three maximal consistent subsets.

The third entails \(q\); the first and second retain \(\neg q\).
The choice criterion of retaining a maximal consistent subset is
met by all three, while the static consequence relation has the
same meaning for all of them. One policy could choose each different
repair according to which statement it treats as least reliable.
Those choices require a reliability ordering or another policy input
not supplied by the consequence relation. This establishes the
claimed underdetermination, not a theorem that every real agent
uses maximal-subset repair or that all three repairs are equally
warranted once further evidence is given.
\end{proof}
\end{annexitem}

\begin{annexitem}{annex:p18.r56.explicit-closure}{thm:p18-r56-explicit-closure}{An explicit-storage limit, not an impossibility of implicit representation}
\begin{proof}
Every valuation making \(p\) true makes \(p\lor q_i\) true,
for each \(i\in\mathbb N\). Hence these are all consequences.
Their expressions are syntactically distinct because their variable
tokens \(q_i\) are distinct. A store with at most \(m\) explicit
formula records cannot contain even the \(m+1\) such records with
indices \(0,\ldots,m\), by the finite pigeonhole principle.
It therefore cannot contain the whole closure.

The premises \(\{p\}\) themselves form a one-record representation.
Moreover a finite instruction which, given a numerical index \(i\),
forms the expression \(p\lor q_i\), describes the displayed
infinite family without storing it as an infinite list. The instruction
only copies \(p\), the disjunction token and the indexed variable
token, so it is a genuinely finite generation procedure. This
distinguishes explicit membership records from a generator or
query procedure. It does not identify the capacities of an actual
person or prove that every logical consequence query is tractable.
\end{proof}
\end{annexitem}

\begin{annexitem}{annex:p18.r56.lexical-representation}{thm:p18-r56-lexical-representation}{The red-shirt inference has explicit formal representations}
\begin{proof}
For the failure without a meaning constraint, take a one-element
first-order domain \(\{u\}\), interpret \(t\) as \(u\),
\(\mathsf{Red}\) as \(\{u\}\), and \(\mathsf{Colored}\)
as empty. The premise is true and the conclusion false, a direct
countermodel in the unrestricted class.

In any interpretation satisfying
\(\forall x(\mathsf{Red}(x)\to\mathsf{Colored}(x))\),
truth of \(\mathsf{Red}(t)\) implies truth of
\(\mathsf{Colored}(t)\) by the universal's instance at the
denotation of \(t\). That universal is satisfied exactly when
the two extensions obey
\(\mathsf{Red}\subseteq\mathsf{Colored}\), giving the equivalent
restricted-model-class statement. This is an explicit first-order
representation with its extra meaning premise or semantic restriction
visible.

For the decomposition version, use
\(\mathsf{Red}(x):=\mathsf{Colored}(x)\land\mathsf{RedHue}(x)\)
as the stated interpretation of redness, not as an equation assumed
for arbitrary independently interpreted predicates. If the conjunction
at \(t\) is true, its first conjunct is true in every interpretation.
Thus the decomposed representation has the ordinary valid
conjunction-elimination form. It provides another actual formalization
once that lexical decomposition is adopted.

These examples suffice to refute an unrestricted impossibility of
formal representation. They do not certify an arbitrary translation
of a real utterance: the intended predicates, context and relevant
lexical relation must still be justified. Nor do these finite
representations prove that all tacit background knowledge admits
this particular kind of decomposition.
\end{proof}
\end{annexitem}

\begin{annexitem}{annex:p18.r56.conditional-repair}{thm:p18-r56-conditional-repair}{Arbitrary deductivization adds the disputed support}
\begin{proof}
Any valuation satisfying every member of \(\Gamma\) makes
\(A_\Gamma\) true, including \(n=0\), when the conjunction
is the truth constant. If it also makes \(A_\Gamma\to\varphi\)
true, material implication forces \(\varphi\) true. This is
exactly the claimed consequence of the enlarged premise set.

If the original inference is invalid, let a countervaluation make
all original premises true and \(\varphi\) false. Its antecedent
\(A_\Gamma\) is true, so it makes \(\chi_\Gamma\) false.
Every original countervaluation has this property. The repair
therefore excludes the countercases by adding a premise whose truth
is precisely in dispute; formal validity alone supplies no reason
for its acceptance.

For a smallest nonempty example, \(p\not\models q\) has the
valuation \(p=1,q=0\). Adding \(p\to q\) yields a valid
inference, but that same valuation falsifies the added premise.
An independently established law, accepted commitment or genuine
lexical relationship can supply an implication in a different case.
Its adequacy is then established separately, not by the bare
observation that adding it makes a derivation valid.
\end{proof}
\end{annexitem}
